Tables such as \tabref{results} can also be used. Be sure to also
refer to the table in the text and offer a discussion about the
results. Note the use of \verb!\tabref{<table label>}! to refefence a
table.

Sometimes it can be nice to include code as well, see
Listing~\ref{CodeTest} for an example. However, think carefully of why
the code should be included. Who will actually read and understand it?

\begin{table}[ht]
  \centering % for centered tables
  \begin{tabular}{c r r r}
    Case & Method \#1 & Method \#2 & Method \#3 \\
    \hline
    1 & 50 & 837 & 970  \\
    2 & 47 & 877 & 230  \\
    3 & 31 & -  & -  \\
    \hline
  \end{tabular}
  \caption{Model results. Note that method \#2 and \#3 did not finish
    for the third case. \blabla{more explanation goes here}
    \DO{Explain all notation in the table, e.g., the dashes in this
      example. You should also refer to the table in the main text.}
    \DONT{Professional tables generally \emph{do not have vertical
        lines}.}}
  \label{tab:results}
\end{table} 

\begin{lstlisting}[language = Matlab, label = CodeTest, caption =
  Sample code.]
function c = repeat(a,b)
%REPEAT Repeat elements of vector.
%   C = REPEAT(A,B) returns A repeated by B, see the example
%   below. Both arguments are assumed to be row vectors of the same
%   length and B must contain nonnegative integers only.
%
%   No error-checking is performed.
%
%   Example:
%     a = [0 1 pi inf nan -nan];
%     b = [1 4 0   2   3    0 ];
%     c = repeat(a,b) % returns c = [0 1 1 1 1 inf inf nan nan nan]

% S. Engblom 2005-10-14

% can't beat this!
ix = cumsum(sparse(1,cumsum([1 b]),1));
c = a(ix(1:end-1));
\end{lstlisting}